\subsection{Exact reduced pentagon assignments in the new quotient}

\paragraph{D23.P7 (a finite calculation with a precise scope).}
Use D23.P6's matrices and put
\[
 X_1=A^2,\quad X_2=B^2,\quad X_3=C^2,\quad
 X_4=(ABA)^2,\quad X_5=(CBC)^2,\qquad J=ABCABA.
\]
Let $F_{\rho,e}$, $\rho\in\mathbb F_3,e\in\mathbb F_2$, be
the six exact first-pair matrices of D23.F13. Let $C_0,Q_0$ be
the full six-dimensional intertwiners of D23.C9. In the finite
derived subgroup $G_{34}'=\langle X_3,X_4\rangle'$, of order
$163296$, consider the equation
\[
 R(JR^{-1}J^{-1})(JF_{\rho,e}^{-1}J^{-1})
 \bigl(Q_0^{-1}rQ_0r^{-1}\bigr)F_{\rho,e}=I,
 \qquad r=C_0^{-1}RC_0.
\tag{D23.P7.1}
\]
For each of the six pairs $(\rho,e)$, exactly $36$ matrices $R$
satisfy this equation. There are no solutions with the right side
replaced by $-I$. Thus the reduced matrix equation admits both
cubic-sign values at every $\rho$.

Every actual multiplier-one full-GT element gives one of these
$216$ assignments, because D23.F13 fixes its first-pair value,
the two-cycle gives the two reversing-half-twist factors, and
D23.C9 gives the displayed fourth factor. This is a necessary
finite assignment theorem. The count by itself does not assert
that every listed assignment comes from a common word, an actual
gentle element, or an actual full-GT element.

\begin{proof}
The reversing half twist satisfies $J^2=-I$ and interchanges
$X_1,X_3$ and $X_4,X_5$ by conjugation while fixing $X_2$.
These are direct matrix identities. Therefore the two-cycle
expresses the five-cycle exactly as in (D23.P7.1), using the
actual inverse $J^{-1}$ rather than silently treating the lift
$J$ as an involution.

The group decomposition D23.C10 gives
\[
 G_{34}'=B_{34}\bigl(H'\times U'\bigr)B_{34}^{-1},
 \qquad |H'|=27,\quad |U'|=6048,
\]
where the two three-dimensional blocks and the splitting matrix
$B_{34}$ are literal matrices in D23.C9's certificate. The
enumeration program reconstructs each of $H',U'$ as the normal
closure of the commutator of its two displayed generators. It
then runs through all $27\cdot6048$ block-diagonal pairs, without
a sampling assumption.

For efficiency the program uses the following equivalent form of
(D23.P7.1). Put
\[
 A_R=RJR^{-1},\qquad
 B_R=rQ_0^{-1}r^{-1}Q_0J.
\]
Multiplying (D23.P7.1) on the indicated sides yields
\[
 F_{\rho,e}A_R=B_RF_{\rho,e}.
\tag{D23.P7.2}
\]
The program tests equality of all entries for all six $F_{\rho,e}$
and every enumerated $R$. It obtains $36$ each; testing the
negative equality gives zero each. Every retained assignment was
then checked again by multiplying the original five factors in
(D23.P7.1), rather than by reusing (D23.P7.2).

All arithmetic is exact over $\mathbb F_3$, after the faithful
restriction-of-scalars embedding of $\mathbb F_9$ matrices. The
reproducible certificate is
\path{agents/pentagon/krammer_pentagon_search.py}; its counts are
in the matching JSON file, and all $216$ matrices with their
$(\rho,e)$ labels are retained in
\path{agents/pentagon/krammer_pentagon_solutions.npz}.

The actual full-GT pentagon lies in the label-five stabilizer
$B_4/\langle z_4\rangle$. The displayed word is a product of
evaluations of $f\in\widehat F_2'$, so its braid exponent sum is
zero. If its lift to $\widehat B_4$ is $z_4^a$, abelianization
gives $12a=0$ in $\widehat{\mathbb Z}$, hence $a=0$. Thus the
literal identity in (D23.P7.1), not merely an arbitrary scalar,
is necessary. In this finite calculation the only additional
possible scalar lift would in any event be $-I$, and it yields
no solutions.
\end{proof}

\paragraph{Unresolved comparison.}
D23.P7 is stronger than discarding the new representation by a
nilpotence criterion: its first-pair cubic sign enters a genuinely
new pentagon calculation. Its outcome nevertheless leaves both
signs. The common-word correlations and the distinction between
finite assignments and actual profinite lifts must still be
respected. In particular no inference $e=\nu_4$ is contained in
the count $36$.
